\textbf{Phần 1. Trắc nghiệm lựa chọn}
\setcounterpageref{ex}{0}

\begin{ex} %[4D11N62]
Kết quả của phép tính $\dfrac{3}{5} + \dfrac{1}{5}$ là
\choice
{$\dfrac{4}{10}$}
{$\dfrac{2}{5}$}
{\True $\dfrac{4}{5}$}
{$\dfrac{3}{10}$}
\loigiai{
$\dfrac{3 + 1}{5} = \dfrac{4}{5}$.\\
\textbf{Đáp án của câu hỏi là C}
}
\end{ex}

\begin{ex} %[4D11N62]
Kết quả của phép trừ $\dfrac{7}{9} - \dfrac{2}{9}$ là
\choice
{$\dfrac{5}{0}$}
{$\dfrac{9}{9}$}
{$\dfrac{4}{9}$}
{\True $\dfrac{5}{9}$}
\loigiai{
$\dfrac{7 - 2}{9} = \dfrac{5}{9}$.\\
\textbf{Đáp án của câu hỏi là D}
}
\end{ex}

\begin{ex} %[4D11H62]
Kết quả của phép tính $\dfrac{1}{2} + \dfrac{3}{4}$ là
\choice
{\True $\dfrac{5}{4}$}
{$\dfrac{4}{6}$}
{$1$}
{$\dfrac{3}{8}$}
\loigiai{
$\dfrac{2}{4} + \dfrac{3}{4} = \dfrac{5}{4}$.\\
\textbf{Đáp án của câu hỏi là A}
}
\end{ex}

\begin{ex} %[4D11H62]
Kết quả của phép tính $1 - \dfrac{4}{7}$ là
\choice
{$\dfrac{4}{7}$}
{$\dfrac{5}{7}$}
{$\dfrac{2}{7}$}
{\True $\dfrac{3}{7}$}
\loigiai{
$\dfrac{7}{7} - \dfrac{4}{7} = \dfrac{3}{7}$.\\
\textbf{Đáp án của câu hỏi là D}
}
\end{ex}

\begin{ex} %[4D11H62]
Tính thuận tiện: $\dfrac{3}{8} + \dfrac{5}{12} + \dfrac{5}{8}$.
\choice
{$1$}
{$\dfrac{13}{12}$}
{\True $\dfrac{17}{12}$}
{$2$}
\loigiai{
$\left(\dfrac{3}{8} + \dfrac{5}{8}\right) + \dfrac{5}{12} = 1 + \dfrac{5}{12} = \dfrac{17}{12}$.\\
\textbf{Đáp án của câu hỏi là C}
}
\end{ex}

\begin{ex} %[4D11H62]
Tìm $x$ biết: $x + \dfrac{1}{3} = \dfrac{5}{6}$.
\choice
{\True $\dfrac{1}{2}$}
{$\dfrac{4}{6}$}
{$\dfrac{2}{3}$}
{$\dfrac{1}{6}$}
\loigiai{
$x = \dfrac{5}{6} - \dfrac{2}{6} = \dfrac{3}{6} = \dfrac{1}{2}$.\\
\textbf{Đáp án của câu hỏi là A}
}
\end{ex}

\begin{ex} %[4D11V62]
Một mảnh vườn hình chữ nhật có chiều dài $\dfrac{4}{5}\text{ m}$, chiều rộng ngắn hơn chiều dài $\dfrac{1}{3}\text{ m}$. Chiều rộng mảnh vườn là
\choice
{$\dfrac{3}{2}\text{ m}$}
{\True $\dfrac{7}{15}\text{ m}$}
{$\dfrac{1}{2}\text{ m}$}
{$\dfrac{2}{15}\text{ m}$}
\loigiai{
$\dfrac{4}{5} - \dfrac{1}{3} = \dfrac{12 - 5}{15} = \dfrac{7}{15}\text{ m}$.\\
\textbf{Đáp án của câu hỏi là B}
}
\end{ex}

\begin{ex} %[4D11V62]
Tính giá trị biểu thức: $\dfrac{7}{10} + \dfrac{1}{2} - \dfrac{3}{5}$.
\choice
{\True $\dfrac{3}{5}$}
{$\dfrac{5}{10}$}
{$\dfrac{4}{5}$}
{$1$}
\loigiai{
$\dfrac{7}{10} + \dfrac{5}{10} - \dfrac{6}{10} = \dfrac{6}{10} = \dfrac{3}{5}$.\\
\textbf{Đáp án của câu hỏi là A}
}
\end{ex}

\begin{ex} %[4D11V62]
Hai tổ công nhân cùng sửa một đoạn đường: Tổ 1 sửa được $\dfrac{3}{8}$ đoạn đường, tổ 2 sửa được $\dfrac{2}{5}$ đoạn đường. Cả hai tổ sửa được số phần đoạn đường là
\choice
{$\dfrac{5}{13}$}
{\True $\dfrac{31}{40}$}
{$\dfrac{29}{40}$}
{$\dfrac{1}{8}$}
\loigiai{
$\dfrac{3}{8} + \dfrac{2}{5} = \dfrac{15 + 16}{40} = \dfrac{31}{40}$.\\
\textbf{Đáp án của câu hỏi là B}
}
\end{ex}

\begin{ex} %[4D11V62]
Một hộp kẹo có 24 cái. Bạn An lấy $\dfrac{1}{3}$ số kẹo, bạn Bình lấy $\dfrac{1}{4}$ số kẹo. Phân số chỉ số kẹo cả hai bạn đã lấy là
\choice
{$\dfrac{2}{7}$}
{$\dfrac{5}{12}$}
{\True $\dfrac{7}{12}$}
{$\dfrac{1}{12}$}
\loigiai{
$\dfrac{1}{3} + \dfrac{1}{4} = \dfrac{7}{12}$.\\
\textbf{Đáp án của câu hỏi là C}
}
\end{ex}

\begin{ex} %[4D11V62]
Tìm $x$ biết: $\dfrac{7}{4} - x = \dfrac{3}{4} + \dfrac{1}{2}$.
\choice
{$\dfrac{3}{4}$}
{$1$}
{$\dfrac{1}{4}$}
{\True $\dfrac{1}{2}$}
\loigiai{
$\dfrac{7}{4} - x = \dfrac{5}{4} \Rightarrow x = \dfrac{7}{4} - \dfrac{5}{4} = \dfrac{2}{4} = \dfrac{1}{2}$.\\
\textbf{Đáp án của câu hỏi là D}
}
\end{ex}

\begin{ex} %[4D11C62]
Tính giá trị biểu thức: $S = \dfrac{1}{2} + \dfrac{1}{6} + \dfrac{1}{12} + \dfrac{1}{20} + \dfrac{1}{30}$.
\choice
{$\dfrac{4}{5}$}
{\True $\dfrac{5}{6}$}
{$\dfrac{6}{7}$}
{$1$}
\loigiai{
$S = \left(1 - \dfrac{1}{2}\right) + \left(\dfrac{1}{2} - \dfrac{1}{3}\right) + \dots + \left(\dfrac{1}{5} - \dfrac{1}{6}\right) = 1 - \dfrac{1}{6} = \dfrac{5}{6}$.\\
\textbf{Đáp án của câu hỏi là B}
}
\end{ex}

\textbf{Phần 2. Trắc nghiệm đúng sai}
\setcounterpageref{ex}{0}

\begin{ex} %[4D11H62]
Xét tính đúng sai của các phép tính:
\choiceTF
{\True $\dfrac{2}{3} + \dfrac{1}{6} = \dfrac{4}{6} + \dfrac{1}{6} = \dfrac{5}{6}$.}
{\True $1 - \dfrac{2}{5} = \dfrac{3}{5}$.}
{$\dfrac{3}{4} - \dfrac{1}{3} = \dfrac{2}{1} = 2$.}
{\True $\dfrac{3}{4} - \dfrac{1}{3} = \dfrac{9 - 4}{12} = \dfrac{5}{12}$.}
\loigiai{
\begin{itemize}
    \item $2/3 + 1/6 = 5/6$. Đúng.
    \item $1 - 2/5 = 3/5$. Đúng.
    \item Phép trừ phân số phải quy đồng, không thể lấy tử trừ tử mẫu trừ mẫu. Ý 3 Sai, Ý 4 Đúng.
\end{itemize}
\begin{itemize}
    \item a) \textbf{Đúng}.
    \item b) \textbf{Đúng}.
    \item c) \textbf{Sai}.
    \item d) \textbf{Đúng}.
\end{itemize}
}
\end{ex}

\begin{ex} %[4D11V62]
Một chai đựng $1\text{ l}$ nước ngọt. Bạn Lan uống $\dfrac{1}{4}\text{ l}$, bạn Mai uống $\dfrac{3}{8}\text{ l}$. Xét tính đúng sai:
\choiceTF
{\True Cả hai bạn uống hết $\dfrac{5}{8}\text{ l}$ nước ngọt.}
{\True Bạn Mai uống nhiều hơn bạn Lan là $\dfrac{1}{8}\text{ l}$.}
{Trong chai còn lại $\dfrac{1}{2}\text{ l}$ nước ngọt.}
{\True Trong chai còn lại $\dfrac{3}{8}\text{ l}$ nước ngọt.}
\loigiai{
\begin{itemize}
    \item Tổng uống: $2/8 + 3/8 = 5/8$ l. Đúng.
    \item Hiệu: $3/8 - 2/8 = 1/8$ l. Đúng.
    \item Còn lại: $1 - 5/8 = 3/8$ l $\neq 1/2$ l. Ý 3 Sai, Ý 4 Đúng.
\end{itemize}
\begin{itemize}
    \item a) \textbf{Đúng}.
    \item b) \textbf{Đúng}.
    \item c) \textbf{Sai}.
    \item d) \textbf{Đúng}.
\end{itemize}
}
\end{ex}

\textbf{Phần 3. Trả lời ngắn}
\setcounterpageref{ex}{0}

\begin{ex} %[4D11H62]
Tính tổng: $\dfrac{1}{3} + \dfrac{2}{5}$ (dạng a/b).
\shortans{11/15}
\loigiai{
\textbf{Đáp số:} $11/15$\\
$5/15 + 6/15 = 11/15$.
}
\end{ex}

\begin{ex} %[4D11H62]
Tính hiệu: $\dfrac{5}{6} - \dfrac{1}{4}$ (dạng a/b).
\shortans{7/12}
\loigiai{
\textbf{Đáp số:} $7/12$\\
$10/12 - 3/12 = 7/12$.
}
\end{ex}

\begin{ex} %[4D11V62]
Tính giá trị: $1 + \dfrac{2}{7} - \dfrac{3}{7}$ (dạng a/b).
\shortans{6/7}
\loigiai{
\textbf{Đáp số:} $6/7$\\
$7/7 + 2/7 - 3/7 = 6/7$.
}
\end{ex}

\begin{ex} %[4D11V62]
Tìm $x$ biết: $x - \dfrac{3}{10} = \dfrac{1}{2}$ (dạng a/b).
\shortans{4/5}
\loigiai{
\textbf{Đáp số:} $4/5$\\
$x = 5/10 + 3/10 = 8/10 = 4/5$.
}
\end{ex}

\begin{ex} %[4D11V62]
Một hình tam giác có chu vi $\dfrac{9}{10}\text{ m}$, hai cạnh có độ dài là $\dfrac{1}{5}\text{ m}$ và $\dfrac{3}{10}\text{ m}$. Tính độ dài cạnh thứ ba (rút gọn tối giản dạng a/b).
\shortans{2/5}
\loigiai{
\textbf{Đáp số:} $2/5$\\
$9/10 - 2/10 - 3/10 = 4/10 = 2/5\text{ m}$.
}
\end{ex}

\begin{ex} %[4D11C62]
Tính nhanh: $\dfrac{1}{3} + \dfrac{1}{4} + \dfrac{2}{3} + \dfrac{3}{4}$.
\shortans{2}
\loigiai{
\textbf{Đáp số:} $2$\\
$(1/3 + 2/3) + (1/4 + 3/4) = 1 + 1 = 2$.
}
\end{ex}
